\section{论文内容与源文件组织}
\label{sec:writing}

\subsection{示例的用途与边界}
毕业论文首先需要提出明确的问题，说明所用材料和方法，并形成可以追溯的论证。模板承担的是呈现这些内容的工作。本示例围绕源文件组织展开，帮助使用者找到题目、正文、图表及文献的修改位置；其中的文字与对象均应在实际写作时替换。本例不是研究报告，不对学习效果、算法性能或管理实践作实证判断。

使用模板时，可以先保留结构，逐步替换内容。主文件列出各部分的先后顺序；配置文件保存题目和作者等信息；正文目录保存各节内容；文献库保存来源的书目信息。这样组织后，修改段落不需要同步改动公共版式。多人协作时，也更容易判断一次修改究竟改变了论证还是格式。

\subsubsection{从问题到材料}
研究问题决定所需证据。以机器学习为学习主题时，可以从已有教材梳理概念，再查阅与具体问题相关的原始文献。《机器学习》与《统计学习方法》提供了不同的教材组织方式，本例引用它们仅用于演示中文图书的脚注及排序。\cuebcite{zhou2016}\cuebcite{li2019} 实际论文应说明每个来源支持哪一项论断，不能把罗列书名当作论证过程。

记录资料时，建议同时保存作者、题名、版本、出版社、年份和已核对的引用页码。页码应来自实际阅读的版本；不同版本之间不可直接照搬。一般说明性内容可以放在普通脚注中。\footnote{本脚注说明模板用法，与文献引用共享脚注计数。它不对应文末的一条书目记录。} 正文的脚注号与文末文献序号承担不同功能，不要求逐项相同。

\subsubsection{从材料到章节}
本例只使用三级标题。一级标题组织核心问题，二级标题划分论证步骤，三级标题解释局部方法。标题需要概括其下内容，不能仅以“相关内容”代替具体主题。尽量避免只有标题而没有正文的层级，也不宜为了版面整齐将每一个自然段都设为独立小节。

\subsection{公式、图与表的配合}
\label{sec:objects}
公式应服务于定义或推导。设人工构造的三个观测值为 $x_1=2$、$x_2=4$、$x_3=6$，其均值仅用于说明数学排版：
\begin{equation}
  \bar{x}=\frac{1}{n}\sum_{i=1}^{n}x_i=4,\qquad n=3.
  \label{eq:mean}
\end{equation}
式\eqref{eq:mean}中的 $n$ 表示观测个数，$x_i$ 表示第 $i$ 个演示值。变量在首次出现时应解释清楚，计量单位和取值范围则应由实际问题确定。这里没有使用真实样本，也没有据此进行统计推断。

图\ref{fig:workflow}用简单的方框表示写作步骤。箭头表达工作顺序，不表示因果效应。该图由本例原创绘制，不依赖外部图片文件。正式论文可以使用矢量图或清晰的位图，同时在正文说明图中信息为何与论证有关。
\begin{figure}[htbp]
  \centering
  \setlength{\fboxsep}{7pt}
  \fbox{确定问题}\enspace$\longrightarrow$\enspace
  \fbox{整理证据}\enspace$\longrightarrow$\enspace
  \fbox{撰写与核对}
  \caption{论文写作步骤示意}
  \label{fig:workflow}
\end{figure}

表\ref{tab:synthetic}展示三线表。读者应能从表题、表头和表注判断数据的含义。本例同时保留样本数与计量单位，避免只显示一列没有解释的数字。图题置于图下，表题置于表上；题名中不应加入手工编号。
\begin{table}[htbp]
  \centering
  \caption{人工构造的演示值}
  \label{tab:synthetic}
  \begin{tabular}{ccc}
    \toprule
    记录编号 & 演示值（单位：点） & 数据来源 \\
    \midrule
    A & 2 & 人工构造 \\
    B & 4 & 人工构造 \\
    C & 6 & 人工构造 \\
    \bottomrule
  \end{tabular}
  \par\smallskip
  {\zihao{5}注：本表仅演示排版，不包含实测数据。}
\end{table}

\subsection{交叉引用与修改过程}
正文通过标签引用对象。插入新的图表以后，编译系统会更新编号，避免“见图2”之类的手写数字滞后。图表仍须靠近首次讨论的位置，但排版系统可能根据页面剩余空间调整其位置。若一个浮动对象出现在下一页，应先检查尺寸和前后内容，谨慎使用强制定位。

术语、符号和数据精度需要全文一致。例如，均值的符号在式\eqref{eq:mean}中已经定义，后文应沿用该符号。不同来源若使用不同定义，应明确区分，不能仅因为名称接近而合并。英文教材同样需要保留准确的书目信息；本例以 Bishop 的图书演示英文作者和书名的著录。\cuebcite{bishop2006}

写作过程中应保存可以恢复的版本，特别是在调整章节结构和文献样式之前。每次修改后先检查相关段落，再编译全文。格式修改可能改变断行和分页，因此局部视觉效果合适并不意味着后续各页仍然正确。第\ref{sec:validation}节给出针对这一问题的检查示例。
